\documentclass[11pt,a4paper]{article}
\usepackage[T1]{fontenc}\usepackage[landscape,margin=12mm]{geometry}
\usepackage{graphicx,hyperref,pgffor}
\setlength{\parindent}{0pt}
\newcommand{\AtlasDataset}[2]{%
  \foreach \exc [count=\groupindex from 1] in {#2}{%
    \edef\groupid{\ifnum\groupindex<10 0\fi\groupindex}%
    \foreach \method in {A,B,C,D,E,F,G,H}{%
      \clearpage
      \section*{Set #1, method \method, excitation \exc\ A}
      \includegraphics[width=\linewidth]{figures/atlas/set#1_group\groupid_\method.png}
      \par\small Original: blue; reconstructed: orange. Errors are signed, in mVs.
      The excitation graph uses measured currents. White regions lack support.
      The d-flux 3D view places q-current on the first axis; heatmaps keep physical d/q order.
    }%
  }%
}
\title{Excitation-Current Coordinate Reduction: Complete Reconstruction Atlas}
\author{Winding Factor Research Study}\date{24 October 2026}
\begin{document}\maketitle
One page per dataset, excitation group and method (A--H). The main paper defines the models, point-based metrics, support rules and independent prediction experiment.
\par\medskip
Both flux components are shown. Left: original (blue) and shared-table reconstruction (orange). Right: signed error over physical stator currents. Units are mVs. Missing regions are not extrapolated. Colour limits are shared across each dataset.
\par\medskip
The figures interpolate each original group's measured excitation alongside its flux. A nominal group label is not substituted for the pointwise measured current. The 41-by-41 display grid is distinct from the 81-by-81 stored reduced grid.
\par\medskip
Axis labels use ``current in A'' and ``flux linkage in mVs''. The d-flux 3D view exchanges the displayed d- and q-current axes to expose the surface; the error maps retain their physical d/q arrangement.
\setkeys{Gin}{height=0.80\textheight,keepaspectratio}
\centering
\AtlasDataset{1}{2,5,8,11,14,17,20,23,26,29,32,36,40}
\AtlasDataset{2}{0,2,5,8,11,14,17,20,23,26,29,32,36,40}
\AtlasDataset{3}{3,6,9,12,15,18,21,24,27,30}

\end{document}
